\chapter{Configuration Generation}

Now that the design option tree has been screened to only allow for plausible subsystem options, some possible configurations can be generated, to be evaluated holistically. If every accepted option in the DOT were mathematically combined, it would result in thousands of potential aircraft layouts, which would be impractical to evaluate. In order to narrow down the subsystems being evaluated, the configurations were built without the architecture-independent features.

Subsystems such as avionics, construction materials, and fuel compatibility were omitted from the configuration generation, as these choices are largely independent of the main system architecture and can be integrated within most configurations.

Therefore, to generate a focused and meaningful set of aircraft configurations, the primary aerodynamic and structural layout drivers were isolated, using the subsystems presented in \autoref{tab:configuration_design_options} 

% \begin{itemize}
%     \item \textbf{Wing Configuration} (Low, Mid, High, Flying Wing)
%     \item \textbf{Empennage Configuration} (T-tail, H-Tail, Conventional, V-Tail, Inverted V-Tail, Canard, No tail)
%     \item \textbf{Propulsion Mounting} (Fuselage Mounted, Buried, Wing Mounted)
%     \item \textbf{Undercarriage} (Skids, Retractable gear)
% \end{itemize}

\begin{table}[H]
    \centering
    \caption{Configuration design options}
    \label{tab:configuration_design_options}
    \begin{tabular}{|l|l|}
        \hline
        \textbf{Design category} & \textbf{Options} \\ \hline \hline
        Wing configuration & Low, Mid, High, Flying Wing \\ \hline
        Empennage configuration & T-tail, H-tail, Conventional, V-tail, Inverted V-tail, Canard, No tail \\ \hline
        Propulsion mounting & Fuselage Mounted, Buried, Wing Mounted \\ \hline
        Undercarriage & Skids, Retractable gear \\ \hline
    \end{tabular}
\end{table}
By systematically iterating through a morphological matrix of these specific parameters, an initial set of 20 distinct aircraft configurations was generated. During this combination process, inherent physical and aerodynamic incompatibilities were immediately filtered out to keep the configurations realistic. For example, a Flying Wing configuration inherently dictates a No Tail criterion. The 20 configurations have been tabulated in \autoref{Aircraft_Configurations_Table}

\section{Shortlisting Configurations}
In order to narrow down the 20 configurations to a shortlist, first, several architectural branches were systematically eliminated due to inherent conflicts with the testbed's primary mission:

\begin{itemize}
    \item \textbf{Mid-Wing Configurations (Configs 14, 15):} Eliminated primarily due to internal packaging constraints. The structural carry-through box of a mid-wing design heavily obstructs the continuous internal volume required for the payload, as stated in  \textbf{HUG-STR-017}. Furthermore, it does not offer distinct aerodynamic advantages over a high or low-wing configuration to justify the harder structural integration and the necessary reinforcements.
    
    \item \textbf{Canard Configurations (Configs 7, 15):} Discarded due to aerodynamic interference. A forward canard generates downwash and a turbulent wake that affects the main wing. Because the core objective of this flying laboratory is to capture clean aeroelastic and flutter data from the main wing, upstream flow disturbances are unacceptable. Additionally, canards are less representative of the future commercial transport architectures this project aims to emulate.
    
    \item \textbf{High-Wing with Fuselage-Mounted Propulsion (Configs 16, 17, 19):} Eliminated due to aerodynamic and structural integration challenges. If the micro-turbine is mounted on the aft fuselage sides, the intakes risk ingesting the turbulent wake generated by the high wing, degrading engine performance. Conversely, a centrally mounted engine placed above the fuselage introduces structural interference with the high wing's carry-through box and modular mounting mechanisms. While physically feasible, these conflicts make the high-wing architecture fundamentally inferior to the low-wing alternatives for this specific propulsion setup.
    
    \item \textbf{Wing-Mounted Propulsion (Configs 9, 18, 20):} Eliminated due to strict versatility and aeroelastic constraints. The drone is required to facilitate the testing of highly flexible, swappable HAR wings. Hard-mounting a jet engine directly to these flexible test articles introduces massive inertial and torsional coupling that would corrupt the baseline flutter data, while simultaneously violating the requirement for wing modularity and versatility.
\end{itemize}


Many of the remaining configurations share fundamental aerodynamic and structural characteristics, allowing them to be grouped into distinct architectural families. To avoid redundant analysis during the detailed trade-off phase, these families were evaluated to identify a few, optimal representative configurations for each. This down-selection process effectively narrowed the field from 20 permutations to the following highly distinct candidate architectures

\begin{itemize}
    \item \textbf{Configuration 1 (Low Wing, T-Tail, Retractable Gear):} Selected to serve as a traditional baseline model. Because this configuration closely mirrors standard transport aircraft architectures, the aeroelastic and flutter data collected will be highly relevant and easily scalable to the commercial airplanes the testbed is designed to emulate. 
    
    \item \textbf{Configuration 3 (Low Wing, H-Tail, Retractable Gear):} Similar to Configuration 1, this serves as a highly relevant baseline model. Retaining the H-tail variant alongside the T-tail allows for further trade-off analysis regarding engine placement and tail-wake interference during the detailed design phase.
    
    \item \textbf{Configuration 8 (Low Wing, Inverted V-Tail, Retractable Gear):} This architecture was selected primarily for its clean upper aerodynamic surface. The inverted V-tail keeps the empennage clear of the upper fuselage, providing an unobstructed region to mount a test airfoil. This allows the experimental wings to be safely tested to structural failure without their wake interfering with the aircraft's pitch and yaw control.
    
    \item \textbf{Configurations 10 \& 11 (Flying Wing, Buried Engine):} Selected for their exceptionally low parasitic drag profile. This high aerodynamic efficiency ensures the aircraft can easily and reliably reach the high Mach numbers necessary to trigger aeroelastic phenomena during flight testing. Both landing gear variants (retractable gear in Config 10 and skids in Config 11) are retained for further evaluation in the mid-term trade-off.
\end{itemize}




\begin{longtable}{|c|c|c|p{2.5cm}|p{4.5cm}|}
    \caption{Preliminary Aircraft Design Configurations} \label{Aircraft_Configurations_Table} \\
    
    \hline
    \textbf{Config. ID} & \textbf{Wing} & \textbf{Propulsion} & \textbf{Empennage} & \textbf{Undercarriage} \\
    \hline\hline
    \endfirsthead

    \hline
    \textbf{Config. ID} & \textbf{Wing} & \textbf{Propulsion} & \textbf{Empennage} & \textbf{Undercarriage} \\
    \hline\hline
    \endhead

    \hline
    \multicolumn{5}{r}{\textit{Continued on next page}} \\
    \endfoot

    \hline
    \endlastfoot

    HUG-CFG-01 & Low wing & Fuselage mounted & T-tail & Retractable gear \\
    \hline
    HUG-CFG-02 & Low wing & Fuselage mounted & T-tail & Skids (Belly landing / Rail launched) \\
    \hline
    HUG-CFG-03 & Low wing & Fuselage mounted & H-tail & Retractable gear \\
    \hline
    HUG-CFG-04 & Low wing & Fuselage mounted & H-tail & Skids (Belly landing / Rail launched) \\
    \hline
    HUG-CFG-05 & Low wing & Fuselage mounted & V-tail & Retractable gear \\
    \hline
    HUG-CFG-06 & Low wing & Fuselage mounted & V-tail & Skids (Belly landing / Rail launched) \\
    \hline
    HUG-CFG-07 & Low wing & Fuselage mounted & Canard & Retractable gear \\
    \hline
    HUG-CFG-08 & Low wing & Fuselage mounted & Inverted V-tail & Retractable gear \\
    \hline
    HUG-CFG-09 & Low wing & Wing Mounted & Conventional tail & Retractable gear \\
    \hline
    HUG-CFG-10 & Flying wing & Buried engine & No tail & Retractable gear \\
    \hline
    HUG-CFG-11 & Flying wing & Buried engine & No tail & Skids (Belly landing / Rail launched) \\
    \hline
    HUG-CFG-12 & Flying wing & Fuselage mounted & No tail & Retractable gear \\
    \hline
    HUG-CFG-13 & Flying wing & Fuselage mounted & No tail & Skids (Belly landing / Rail launched) \\
    \hline
    HUG-CFG-14 & Mid wing & Fuselage mounted & H-tail & Retractable gear \\
    \hline
    HUG-CFG-15 & Mid wing & Fuselage mounted & Canard & Skids (Belly landing / Rail launched) \\
    \hline
    HUG-CFG-16 & High wing & Fuselage mounted & V-tail & Skids (Belly landing / Rail launched) \\
    \hline
    HUG-CFG-17 & High wing & Fuselage mounted & H-tail & Skids (Belly landing / Rail launched) \\
    \hline
    HUG-CFG-18 & High wing & Wing Mounted & T-tail & Retractable gear \\
    \hline
    HUG-CFG-19 & High wing & Fuselage mounted & Inverted V-tail & Skids (Belly landing / Rail launched) \\
    \hline
    HUG-CFG-20 & High wing & Wing Mounted & Conventional tail & Retractable \\
    \hline

\end{longtable}